\chapter{Алгоритм построения Парето-границы}\label{ch:algorithm}

В предыдущей главе доказана структурная Теорема~\ref{thm:main} о форме
оптимального waveform: оптимум в классе $\varepsilon$-зарядово-сбалансированных
управлений релейный с числом активных фаз
$K_{\mathrm{eff}} \le d_{\mathrm{state}} + 2 = 4$. Это сужает пространство
поиска от континуума до конечного множества кандидатов, так что полная
Парето-граница задачи многокритериальной оптимизации обновления EPD-панели
строится конструктивно.

Глава содержит постановку этой задачи (\S\ref{sec:moo-formulation}), описание
метода $\varepsilon$-ограничений
(\S\ref{sec:eps-constraint})~\cite{vasiliev2002_methods}, анализ свойств
получаемого фронта (\S\ref{sec:pareto-properties}), псевдокод алгоритма с
оценкой сложности (\S\ref{sec:algorithm-complexity}) и сравнение с
генетическим алгоритмом NSGA-II Деба и
соавторов~\cite{deb2002_nsga2} (\S\ref{sec:nsga2-comparison}).


\section{Многокритериальная постановка}\label{sec:moo-formulation}

Используя обозначения из~\ref{sec:theorem-formulation}, обозначим через
$u = \bigl((V[k], T[k])\bigr)_{k=0}^{K-1}$ дискретный waveform длины $K$,
где $V[k] \in \mathcal{U}$ и $T[k] \in \mathcal{T}_{\mathrm{disc}}$. Для
каждой пары начального и целевого состояний пикселя $(z_{\mathrm{init}},
z_{\mathrm{target}})$ задача формулируется как:
\begin{equation}\label{eq:moo}
    \min_{u \in \mathcal{U}_{\mathrm{feas}}}
        \bigl( E(u),\; G(u; z_{\mathrm{init}}, z_{\mathrm{target}}),\; \tau(u) \bigr),
\end{equation}
где $E(u)$~--- суммарная энергия из~\eqref{eq:phase-energy},
$G(u; \cdot)$~--- метрика гостинга из~\eqref{eq:ghost-cost} и $\tau(u)$~---
суммарная длительность переходного процесса из~\eqref{eq:latency}.
Множество допустимых управлений $\mathcal{U}_{\mathrm{feas}}$ задаётся
совокупностью ограничений из~\eqref{eq:optimization-problem}: дискретность
$V[k]$ и $T[k]$, $\varepsilon$-зарядовый баланс
из~\eqref{eq:hard-constraint} и оценка числа активных фаз
из Теоремы~\ref{thm:main}.

\begin{definition}[Парето-доминирование]\label{def:pareto-dom}
    Управление $u_1$ \emph{доминирует} $u_2$ (записывается $u_1 \prec u_2$),
    если все три критерия~\eqref{eq:moo} удовлетворяют $f_i(u_1) \le f_i(u_2)$
    при $i=1,2,3$ и хотя бы одно неравенство строгое. Множество
    \emph{Парето-оптимальных} управлений определяется как
    \begin{equation*}
        \mathcal{P}^{*} = \bigl\{ u \in \mathcal{U}_{\mathrm{feas}} :
            \nexists\, u' \in \mathcal{U}_{\mathrm{feas}} \text{ такое, что } u' \prec u
        \bigr\}.
    \end{equation*}
\end{definition}

Образ множества $\mathcal{P}^{*}$ в пространстве критериев
$\bigl(E, G, \tau\bigr)$ называется \emph{Парето-границей}; именно этот
геометрический объект и требуется построить.


\section{Метод $\varepsilon$-ограничений}\label{sec:eps-constraint}

Классический скаляризующий приём (см. Васильев~\cite{vasiliev2002_methods},
гл.~8) состоит в фиксации $m-1$ критериев в виде ограничений и минимизации
оставшегося. В трёхкритериальной задаче~\eqref{eq:moo} выберем энергию~$E$ как
минимизируемую цель, а гостинг~$G$ и длительность~$\tau$ переведём в
ограничения:
\begin{equation}\label{eq:eps-constraint-problem}
    \min_{u \in \mathcal{U}_{\mathrm{feas}}} E(u),
    \qquad
    \text{при } G(u; z_{\mathrm{init}}, z_{\mathrm{target}}) \le \varepsilon_G,
    \quad \tau(u) \le \varepsilon_{\tau}.
\end{equation}

Перебирая $(\varepsilon_G, \varepsilon_{\tau})$ по двумерной сетке
$\mathcal{G}_{\varepsilon}$, получаем семейство задач, каждая из которых даёт
одну точку Парето-границы. Число узлов
$|\mathcal{G}_{\varepsilon}| = N_G \cdot N_{\tau}$ выбирается из компромисса
между плотностью фронта и вычислительными затратами (в расчётах
главы~\ref{ch:simulation} используется сетка порядка $8 \times 8$ узлов).

Здесь важно, что каждая частная
задача~\eqref{eq:eps-constraint-problem} при фиксированных
$(\varepsilon_G, \varepsilon_{\tau})$ совпадает с постановкой дискретной
задачи оптимального управления из~\eqref{eq:optimization-problem} с весовым
вектором $(\alpha, \beta, \gamma) = (1, 0, 0)$ при добавочных ограничениях на
$G$ и~$\tau$. Поэтому к ней применима Теорема~\ref{thm:main}: оптимум релейный
с $K_{\mathrm{eff}} \le 4$, и перебор кандидатов ограничен сверху величиной
$|\mathcal{U}|^{K_{\mathrm{eff}}} \cdot
|\mathcal{T}_{\mathrm{disc}}|^{K_{\mathrm{eff}}}$ (а с фильтром
$\varepsilon$-баланса~--- ещё на порядок меньше), так что каждая частная
задача решается перебором за полиномиальное время, а вся Парето-граница~---
за время, линейное по числу узлов сетки.


\section{Свойства Парето-границы}\label{sec:pareto-properties}

Построенная алгоритмом Парето-граница обладает несколькими полезными для
анализа структурными свойствами.

\paragraph{Свойство~1 (монотонность по числу активных фаз).}
Если оптимум на $(\varepsilon_G, \varepsilon_{\tau})$-узле достигается при
$K_{\mathrm{eff}}^{*} = k$, то ослабление любого из ограничений
($\varepsilon_G \uparrow$ или $\varepsilon_{\tau} \uparrow$) не увеличивает
$K_{\mathrm{eff}}^{*}$ нового оптимума; это прямое следствие включения множеств
допустимых waveform'ов: при более слабых ограничениях допустимы все
waveform'ы со старым $K_{\mathrm{eff}}^{*}$, а среди них уже есть оптимальный
по минимуму $E$.

\paragraph{Свойство~2 (нижняя оценка через $G^{*}$).}
Из Теоремы~\ref{thm:lower-bound} в любой точке фронта выполнено
$E^{*} \ge K_{\mathrm{lower}} \cdot G^{*}$ с положительной константой
$K_{\mathrm{lower}}$, оцениваемой по измеренной на стенде зависимости
$E^{*}(C^{*})$ (\S\ref{sec:result}). Это даёт универсальный <<пол>> кривой
$E^{*}(G^{*})$: ни одна точка фронта не может лежать ниже прямой
$E = K_{\mathrm{lower}} \cdot G$.

\paragraph{Свойство~3 (доминирование заводского LUT).}
Если найденная Парето-граница $\mathcal{P}^{*}_{\mathrm{ours}}$ строго
доминирует точку, соответствующую заводскому LUT~$B_0$, то любой выбор
$(\alpha, \beta, \gamma) \ne 0$ в~\eqref{eq:cost} даёт энергетический
выигрыш. Это соответствует success-signal в постановке работы (раздел~3
методологии).


\section{Псевдокод и оценка сложности}\label{sec:algorithm-complexity}

Конструктивный алгоритм представлен блок-схемой на
Рис.~\ref{fig:algorithm-flowchart}. Вход: пара
$(z_{\mathrm{init}}, z_{\mathrm{target}})$, сетка $\mathcal{G}_{\varepsilon}$,
surrogate-модель $\widehat{F}$ (см. \S\ref{sec:reduction}) и параметры
дискретизации $(\mathcal{U}, \mathcal{T}_{\mathrm{disc}})$. Выход: множество
точек Парето-фронта вместе с породившими их waveform'ами.

\begin{figure}[ht]
    \centering
    \begin{tikzpicture}[
        node distance=4mm,
        every node/.style={font=\footnotesize\sffamily},
        box/.style={rectangle, draw, rounded corners=2pt, minimum width=42mm,
            minimum height=6mm, align=center, inner sep=1.5pt},
        arrow/.style={-{Stealth[length=2mm]}, thick},
    ]
    \node[box] (input) {Вход: $(z_{\mathrm{init}}, z_{\mathrm{target}})$, $\mathcal{G}_{\varepsilon}$, $\widehat{F}$};
    \node[box, below=of input] (epsloop) {Цикл по узлам $(\varepsilon_G, \varepsilon_{\tau}) \in \mathcal{G}_{\varepsilon}$};
    \node[box, below=of epsloop] (kloop) {Перебор $K_{\mathrm{eff}} \in \{1,2,3,4\}$};
    \node[box, below=of kloop] (gen) {Генерация bang-bang кандидатов $(V_{1..K}, T_{1..K})$};
    \node[box, below=of gen] (filter) {Фильтр $|\sum V_k T_k| \le \varepsilon$ и $\tau \le \varepsilon_{\tau}$};
    \node[box, below=of filter] (eval) {Оценка $\widehat{F}(u) \to (E, G, \tau)$};
    \node[box, below=of eval] (minE) {Обновление $\arg\min E$ при $G \le \varepsilon_G$};
    \node[box, below=of minE] (output) {Выход: $\mathcal{P}^{*}$ --- точки Парето};
    \draw[arrow] (input) -- (epsloop);
    \draw[arrow] (epsloop) -- (kloop);
    \draw[arrow] (kloop) -- (gen);
    \draw[arrow] (gen) -- (filter);
    \draw[arrow] (filter) -- (eval);
    \draw[arrow] (eval) -- (minE);
    \draw[arrow] (minE) -- (output);
    \end{tikzpicture}
    \caption{Блок-схема алгоритма построения Парето-границы
        методом $\varepsilon$-ограничений с использованием
        Теоремы~\ref{thm:main}}
    \label{fig:algorithm-flowchart}
\end{figure}

\paragraph{Оценка вычислительной сложности.}
Для одного узла $(\varepsilon_G, \varepsilon_{\tau})$ число кандидатов
ограничено сверху как $|\mathcal{U}|^{K_{\mathrm{eff}}} \cdot
|\mathcal{T}_{\mathrm{disc}}|^{K_{\mathrm{eff}}}$. При типовых параметрах
$|\mathcal{U}| = 5$, $|\mathcal{T}_{\mathrm{disc}}| = 16$,
$K_{\mathrm{eff}} = 4$ это даёт $\le 5^4 \cdot 16^4 \approx 4 \cdot 10^7$
кандидатов. Фильтр $\varepsilon$-баланса отсеивает $\approx 80\%$, фильтр
по~$\tau$~--- ещё значительную долю, так что фактическое число вызовов
$\widehat{F}$ составляет $\sim 3 \cdot 10^3$ на узел. На сетке порядка
$8 \times 8$ набирается около $2 \cdot 10^5$ вызовов, что при дешёвой модели $\widehat{F}$
даёт суммарное время менее секунды на пару
$(z_{\mathrm{init}}, z_{\mathrm{target}})$. Это оценка для \emph{общего}
многофазного случая; в калиброванной реализации главы~\ref{ch:simulation}
пространство дополнительно редуцируется до двух параметров, и поиск ещё дешевле.

В символьной форме оценка записывается как $O\bigl(N_G \cdot N_{\tau} \cdot
|\mathcal{U}|^{K_{\mathrm{eff}}} \cdot |\mathcal{T}_{\mathrm{disc}}|^{K_{\mathrm{eff}}}
\cdot C_{\widehat{F}}\bigr)$, где $C_{\widehat{F}}$~--- стоимость одного
вызова surrogate. Благодаря Теореме~\ref{thm:main} показатель степени
$K_{\mathrm{eff}}$ не зависит от длины waveform, что отличает подход от
полного перебора по длине LUT.


\section{Сравнение с генетическим алгоритмом NSGA-II}\label{sec:nsga2-comparison}

В качестве \emph{внешнего baseline} выбран классический эволюционный
многокритериальный алгоритм NSGA-II Деба и
соавторов~\cite{deb2002_nsga2}, не использующий априорной информации о
структуре оптимума. Сравнение приведено в
Таблице~\ref{tab:algorithm-complexity}.

\begin{table}[ht]
    \centering
    \caption{Сравнение нашего $\varepsilon$-constraint
        алгоритма с NSGA-II}\label{tab:algorithm-complexity}
    \footnotesize
    \begin{tabularx}{\textwidth}{lXX}
        \hline
        Параметр & $\varepsilon$-constraint + PMP (наш) & NSGA-II~\cite{deb2002_nsga2} \\
        \hline
        Использование структуры & Да (Теорема~\ref{thm:main}) & Нет \\
        Гарантия глобальной оптимальности & Да (конечный перебор) & Нет (стохастика) \\
        Сложность & Полиномиальная по
            $|\mathcal{U}|, |\mathcal{T}_{\mathrm{disc}}|, K_{\mathrm{eff}}$ &
            $O(M \cdot N_{\mathrm{pop}}^2 \cdot N_{\mathrm{gen}})$ \\
        Воспроизводимость & Детерминированна & Зависит от random seed \\
        Адаптация к новому функционалу & Не нужна (структура та же) &
            Требуется новый запуск \\
        Память & $O(|\mathcal{P}^{*}|)$ & $O(N_{\mathrm{pop}})$ \\
        \hline
    \end{tabularx}
\end{table}

Теоретически NSGA-II уступает структурно ориентированному алгоритму:
эволюционная схема <<вслепую>> исследует тот же конечный набор кандидатов,
который у нас перечисляется явно. Структурная Теорема~\ref{thm:main} сводит
пространство поиска к нескольким целочисленным параметрам
(\S\ref{sec:reduction}), так что исчерпывающий перебор оказывается дешёвым и
гарантирует глобальный оптимум. Стохастические метаэвристики при такой
размерности избыточны и дают лишь приближение без гарантий, так что NSGA-II
остаётся концептуальным baseline, а не практически необходимым инструментом.
